% Source: https://artofproblemsolving.com/community/c3h1771440_2018_school_team
% Stable identity: aops-topic-1771440-q01
Call the tables $T_1,T_2,T_3$ with $1,2,3$ students. Terri joins them with probabilities $\dfrac{1}{6},\dfrac{2}{6},\dfrac{3}{6}$.
\begin{itemize}
\item Terri joins $T_3$ (probability $\dfrac{1}{2}$): $T_3$ already has $4$ students.
\item Terri joins $T_2$ (probability $\dfrac{1}{3}$): the tables have $1,3,3$ students out of $7$, and we need Ursula to join one of the tables with $3$: probability $\dfrac{6}{7}$.
\item Terri joins $T_1$ (probability $\dfrac{1}{6}$): the tables have $2,2,3$, and Ursula must join the table with $3$: probability $\dfrac{3}{7}$.
\end{itemize}
The total probability is
\[
\dfrac{1}{2}+\dfrac{1}{3}\cdot\dfrac{6}{7}+\dfrac{1}{6}\cdot\dfrac{3}{7}=\dfrac{7}{14}+\dfrac{4}{14}+\dfrac{1}{14}=\boxed{\dfrac{6}{7}}.
\]
(Complementary check: no table reaches $4$ only if Terri sits at $T_2$ and Ursula at $T_1$, or Terri at $T_1$ and Ursula at $T_1$ or $T_2$: $\dfrac{1}{3}\cdot\dfrac{1}{7}+\dfrac{1}{6}\cdot\dfrac{4}{7}=\dfrac{1}{21}+\dfrac{2}{21}=\dfrac{1}{7}$.)
